\begin{solapartat}
\punts{0,75} Sabem que el potencial creat per una càrrega puntual s'expressa així:
\[ V = k\,\frac{Q}{r} \]
En aquesta fórmula, la constant de Coulomb ($k$) i la distància ($r$) són magnituds positives. Per tant, el signe del potencial depèn del signe de la càrrega. Així, a prop d'una càrrega positiva el potencial ha de ser positiu i a prop d'una càrrega negativa el potencial ha de ser negatiu. Així doncs, $Q_1$ és una càrrega positiva i $Q_2$ és una càrrega negativa.

\destaca{Alternativament}

Sabem que el camp elèctric apunta en la direcció que disminueix el potencial i podem veure que quan ens movem de $Q_1$ cap a $Q_2$ el potencial disminueix. Per tant, el camp es dirigeix de $Q_1$ cap a $Q_2$. Per una altra banda, les línies de camp surten de les càrregues positives i van a parar a les negatives; per tant, $Q_1$ és una càrrega positiva i $Q_2$ és una càrrega negativa.

\punts{0,25} Per una banda, el camp elèctric ha de ser perpendicular a les línies equipotencials. Al punt $A$ les línies equipotencials són aproximadament verticals. Per tant, al punt $A$ el camp elèctric ha de ser aproximadament horitzontal. En la correcció s'ha de valorar que el camp elèctric s'hagi representat clarament perpendicular a les línies equipotencials.

\punts{0,25} Per una altra banda, ha d'apuntar en la direcció en què el potencial disminueix (o de les càrregues positives cap a les negatives). Per tant, el sentit és el que s'indica a la figura següent:
\figura[0.55]{sol_fig1.png}
\end{solapartat}

\begin{solapartat}
\punts{0,75} El treball fet pel camp elèctric s'expressa així:
\begin{align*}
W_{A\to B} &= -q\Delta V = -(-e)(V_B - V_A) = 1,602\times 10^{-19}\,(10 - (-10))\\
&= 3,204\times 10^{-18}\ \mathrm{J}
\end{align*}

\punts{0,5} El treball és zero perquè la diferència de potencial entre el punt inicial i el final és zero.
\end{solapartat}
